\section{Conclusion}
\label{sec:conclusion}

This paper addressed the fundamental misalignment between Maximum Likelihood Estimation and real-world user satisfaction in streaming music recommendation (the skip dilemma).

Our primary scientific and empirical conclusions are:
\begin{enumerate}
    \item \textbf{Achieving Zero Alignment Tax:} Direct Preference Optimization with an auxiliary SFT anchor ($\gamma=0.10$) proves that optimizing for listening duration does not require sacrificing retrieval precision. On 9,052 held-out test users of Yambda-50M, DPO delivered statistically significant gains in ranking order metrics (\textbf{+0.0017 NDCG@10}, $p = 0.0366$; \textbf{+0.0019 MRR}, $p = 0.0260$) alongside next-item retrieval gains (19.37\% $\to$ 19.52\% HR@10, net $+14$ users via McNemar's test).
    \item \textbf{DPO Superiority over Two-Stage Paradigms:} End-to-end policy tuning over all 43.2M parameters outperformed both Critic Re-Ranking and Actor-Critic PPO, eliminating reward hacking and bypassing frozen embedding bottlenecks.
    \item \textbf{Catalog Health Expansion:} Contrary to concerns of popularity collapse, DPO reduced Spearman correlation with popularity by \textbf{$-13.4\%$} while uncovering \textbf{+197 unique tracks} from the catalog long tail~\cite{petrov2024aligning}.
    \item \textbf{Autoregressive Robustness:} 5-step unrolls demonstrated that the policy does not suffer from compounding exposure bias, delivering a \textbf{+2.1-second} increase in cumulative session playtime per 5 tracks~\cite{matveev2026session}.
\end{enumerate}

\subsection{Engineering Recommendations for Production}
For production recommendation teams, we recommend:
\begin{itemize}
    \item Applying margin-threshold filtering ($\Delta \ge 0.40$) to discard noisy, ambiguous preference pairs.
    \item Maintaining an auxiliary SFT anchor ($\gamma \in [0.05, 0.10]$) to protect policy language capabilities.
    \item Favoring end-to-end direct preference alignment over two-stage frozen-embedding critics.
\end{itemize}

\subsection{Future Directions}
Promising future avenues include incorporating multimodal acoustic embeddings to address cold-start tracks, and deploying online contextual bandits for real-time streaming alignment.
